On considère $f(x)=-x^3+x^2+x+1$. Écris \texttt{dichotomie(debut, fin, precision)} pour renvoyer un encadrement d'une racine, de largeur au plus \texttt{precision}. On suppose \texttt{debut < fin}, \texttt{precision > 0}, $f$ continue et strictement monotone sur l'intervalle, et $f(\texttt{debut})f(\texttt{fin})<0$. Si un milieu est une racine exacte, renvoie ce milieu deux fois.

La précision demandée doit rester atteignable avec les nombres flottants utilisés.
